\section{Phase 1 Strategy}
\label{sec:phase1}

Phase~1 strategy in German Whist is fundamentally about \emph{hand
construction}---assembling the strongest possible 13-card hand for Phase~2.
Since Phase~1 tricks don't count, winning or losing a trick matters only
insofar as it determines which card you receive.

\subsection{The Win/Lose Decision}

Each Phase~1 trick presents a binary strategic choice: should you try to win
(and take the known face-up card) or lose (and take the unknown face-down
card)?

\subsubsection{When to Win}

The face-up card should be won when it is:
\begin{itemize}
  \item \textbf{High-value:} Aces and Kings are almost always worth winning,
    as they are likely trick-winners in Phase~2.
  \item \textbf{Trump:} Any trump card adds to trump control, the strongest
    predictor of Phase~2 success.
  \item \textbf{Synergistic:} Cards that extend existing long suits or fill
    gaps in strong suits are more valuable than isolated high cards.
\end{itemize}

\subsubsection{When to Lose}

Deliberately losing is correct when:
\begin{itemize}
  \item The face-up card is low and in a suit you don't want to extend.
  \item Winning would require expending a valuable card that's better saved
    for Phase~2 or a later Phase~1 trick.
  \item The face-down card has positive expected value (which it does when
    the average remaining unknown card is better than the face-up card).
\end{itemize}

\subsection{Card Valuation}

Not all cards are equally valuable. Our simulation data reveals a hierarchy:

\begin{center}
\begin{tabular}{lc}
\toprule
Card Category & Relative Value \\
\midrule
Aces (trump) & 10.0 \\
Aces (non-trump) & 6.7 \\
Kings (trump) & 7.0 \\
Kings (non-trump) & 4.7 \\
Queens, Jacks & 2.0--3.3 \\
Mid-range (7--10) & 0.5--1.5 \\
Low cards (2--6) & 0.3--0.5 \\
\bottomrule
\end{tabular}
\end{center}

Critically, the value of a card depends on context. A 2 of trumps can be
more valuable than a Queen of a side suit if it enables ruffing.

\subsection{Information Management}

A subtle aspect of Phase~1 is \emph{information asymmetry}. When you win a
trick, both players know what card you received (the face-up card). When you
lose, only you know what you received (the face-down card). This means:
\begin{itemize}
  \item Losing tricks selectively conceals information about your hand.
  \item Winning tricks gives the opponent information but gives you a better
    card (on average, the known face-up card is more desirable than the
    unknown face-down card because you can choose when to win).
\end{itemize}

\subsection{Lead Control}

Maintaining the lead in Phase~1 provides two advantages:
\begin{enumerate}
  \item You see the face-up card before deciding whether to try to win.
  \item You choose which suit to play, potentially forcing the opponent to
    play in suits where they're weak.
\end{enumerate}

However, leading is not always advantageous. If the face-up card is
undesirable, being the follower lets you control whether you win without
having committed a card first.
